% !TeX spellcheck = it_IT

\documentclass[handout]{beamer}

\usepackage[utf8]{inputenc}
\usepackage[italian]{babel}
\usepackage{multicol}
\usetheme{metropolis}
\usepackage{forest}

\usepackage{listings}
\usepackage{color} %red, green, blue, yellow, cyan, magenta, black, white
\definecolor{mygreen}{RGB}{28,172,0} % color values Red, Green, Blue
\definecolor{mylilas}{RGB}{170,55,241}
\lstset{language=Matlab,%
	%basicstyle=\color{red},
	breaklines=true,%
	morekeywords={matlab2tikz},
	keywordstyle=\color{blue},%
	morekeywords=[2]{1}, keywordstyle=[2]{\color{black}},
	identifierstyle=\color{black},%
	stringstyle=\color{mylilas},
	commentstyle=\color{mygreen},%
	showstringspaces=false,%without this there will be a symbol in the places where there is a space
	numbers=left,%
	numberstyle={\tiny \color{black}},% size of the numbers
	numbersep=9pt, % this defines how far the numbers are from the text
	emph=[1]{for,end,break},emphstyle=[1]\color{red}, %some words to emphasise
	%emph=[2]{word1,word2}, emphstyle=[2]{style},
	tabsize=2,
}



\usepackage{tikz}
\usetikzlibrary{shapes,arrows,positioning}

\tikzstyle{startstop} = [ellipse, draw, fill=bg, text=structure, text width=3em, text badly centered, inner sep=0pt]
\tikzstyle{block} = [rectangle, draw, fill=bg, text=structure, text width=5em, text centered]
\tikzstyle{decision} = [diamond, aspect=2, draw, fill=bg, text=structure, text width=3em, text badly centered, inner sep=0pt]
\tikzstyle{input} = [rectangle, draw, fill=bg, text=structure, text width=4em, text centered, rounded rectangle, rounded rectangle east arc=none]
\tikzstyle{output} = [rectangle, draw, fill=bg, text=structure, text width=4em, text centered, rounded rectangle, rounded rectangle west arc=none]

\tikzstyle{line} = [draw, -latex']


\title{Informatica B - Esercitazione 11}
\subtitle{Riepilogo}

\author{Stefano Cereda\\
		stefano1.cereda@mail.polimi.it
	}
\date{19/12/2017}
\institute[PoliMi]{\vspace{0.5cm}\centering Politecnico di Milano \\ \vspace{0.2cm}
	\includegraphics[width=0.2\linewidth]{./logopolimi}}
%\logo{\includegraphics[width=15mm]{./logopolimi}}

\setbeamercovered{invisible}

\makeindex

\begin{document}
\begin{frame}
	\maketitle
\end{frame}

\begin{frame}{Matlab e binario}
\framebreak
Scrivere  una  funzione  Matlab  controllaCP2  che  prende  in  ingresso  una  stringa  (ad  
esempio ‘0101001’)  corrispondente  ad  un  numero  codificato  in  CP2  e  restituisce  due  
valori: 
\begin{enumerate}
	\item par  che  vale  true  se  il  numero  è  pari  e  false  altrimenti  
	\item    pos  che  vale  true  se  il  numero  è  positivo  e  false  altrimenti  
\end{enumerate}
(L’esecuzione  della  funzione  nel  caso  dell’esempio  darebbe  par = false  e  pos =
true)


\framebreak

La funzione matlab deve semplicemente controllare il bit a sinistra per sapere se il numero è positivo, ed il bit a destra per sapere se è pari (ragionate sul perché un numero positivo pari mantiene il primo bit=0 quando passo al negativo).
\begin{lstlisting}
function [par, pos] = controllaCP2(stringa)
if stringa(1) == '0'
pos = true;
else
pos = false;
end
% versione equivalente:
% pos = stringa(1) == '0';

par = stringa(end) == '0';
end
\end{lstlisting}


\end{frame}



\begin{frame}[allowframebreaks]{Matlab - Ricorsione}
	Un metodo per calcolare il valore approssimato della radice quadrata di un
	numero reale non negativo z, detto metodo babilonese, utilizza la seguente
	relazione ricorsiva:
	\[x(n) = \frac{1}{2}(x(n-1)+\frac{z}{x(n-1)})\]
	\[x(0)=1\]
	
	Dove $n$ è un intero non-negativo arbitrario, mentre $x(n)$ rappresenta un'approssimazione della radice quadrata del numero $z$. Per esempio, se $z$ è pari a 2 e $n$ è pari a 3 il valore di $x(n)$ ottenuto applicando la formula indicata sopra è 1.4142 (che è una buona approssimazione della radice di 2).
	
	\framebreak
	
	L'errore di approssimazione associato a $x(n)$ è definito come segue:
	\[errore(n) = | x(n) - x(n-1)|\]
	e decresce al crescere di $n$. Per esempio, se si calcola la radice di 2 con $n$ pari a 1 si ottiene il valore 1.5000 con errore 0.5000. Se invece si calcola la radice dello stesso valore con $n$ pari a 2 si ottiene il valore 1.4167 con errore 0.0833.
	
	\framebreak
	
	\begin{enumerate}
		\item Si scriva una funzione ricorsiva \emph{sqrt1(z,n)} per MATLAB/Octave che
		restituisca un'approssimazione della radice quadrata di $z$ e il relativo errore utilizzando il metodo babilonese. Quando $n$ è pari a zero, si assuma un valore dell’errore pari a inf.
		\item Si scriva una funzione \emph{sqrt2(z,err)} per MATLAB/Octave che, utilizzando la funzione \emph{sqrt1}, restituisca un valore approssimato della radice quadrata di $z$ con errore non superiore al valore $err$ fornito come parametro.
	\end{enumerate}
	
	Nella soluzione di entrambi i quesiti non è permesso utilizzare alcuna funzione di
	libreria di MATLAB/Octave a parte la funzione \emph{abs(x)} per il calcolo del valore
	assoluto di x.
\end{frame}

\begin{frame}{Matlab - Grafici}
	Scrivere una funzione Octave con parametri $a$, $b$ e $p$ che esegue le seguenti operazioni:
	\begin{enumerate}
	\item  calcola i valori $\cos(x^2)$ per valori di x compresi fra a e b, con passo p, e li memorizza in un vettore y
	\item visualizza su un grafico l'andamento di y;
	\item restituisce al chiamante il massimo valore di $\cos(x^2)$ per $a \leq x \leq b$ ed il valore di x corrispondente. Per effettuare il
	calcolo in questione si può utilizzare la funzione max che, quando prende come parametro un vettore, restituisce due
	valori di cui il primo è il valore massimo trovato nel vettore ed il secondo è l’indice della posizione del vettore in cui tale
	valore è contenuto.
\end{enumerate}
\end{frame}


\begin{frame}{Matlab - Scomposizione in fattori}
	Scrivere una funzione che scomponga in fattori primi un numero.
	Il risultato dovrà composto da due vettori: il primo conterrà le basi dei fattori, il secondo conterrà gli esponenti dei fattori, cioè quante volte essi sono ripetuti.
	
	Esempio: il numero 100 è uguale a $2^2 * 5^2$, per cui il vettore delle basi sarà [2 5] mentre quello degli esponenti corrispondenti sarà [2 2].
	
	Un modo semplice per fattorizzare un numero consiste nel dividerlo ripetutamente per tutti i numeri che gli sono inferiori fino ad ottenere 1. 
\end{frame}

\end{document}